%% ECON 282E -- Session 3, Deck A4: Worked Example -- the RBC Model, Fully Discretized
%% Added 2026-09-17 at the instructor's request: B3 referred to an RBC example
%% without ever showing it.  This deck IS the example, line by line.
%% Frame budget deliberately exceeded across Session 3: content has priority and
%% the instructor decides coverage in the room.
%% Calibration: A=1, alpha=0.36, beta=0.95, delta=0.10 -- note delta < 1, unlike
%% 3.A1-3.A3, so there is a closed-form STEADY STATE but no closed-form policy.
%% Every number from tools/figures/s03_numbers.json keys "rbc", "rbc_stochastic".

%% ECON 282E -- Foundations of Macroeconomics (UC Riverside, Fall 2026)
%% Shared Beamer preamble for every deck in slides/S01 ... slides/S10.
%%
%% Usage (a deck lives in slides/S0X/ and is compiled from that folder):
%%   %% ECON 282E -- Foundations of Macroeconomics (UC Riverside, Fall 2026)
%% Shared Beamer preamble for every deck in slides/S01 ... slides/S10.
%%
%% Usage (a deck lives in slides/S0X/ and is compiled from that folder):
%%   %% ECON 282E -- Foundations of Macroeconomics (UC Riverside, Fall 2026)
%% Shared Beamer preamble for every deck in slides/S01 ... slides/S10.
%%
%% Usage (a deck lives in slides/S0X/ and is compiled from that folder):
%%   \input{../preamble/course_preamble.tex}
%%   \decktitle{1}{A1}{Who I Am \& Why This Course}{September 24, 2026}
%%   \begin{document}
%%   \makedecktitle
%%   ...
%%
%% Adapted from Courses/AI-ECON-2026/source/shared_preamble.tex (Metropolis theme,
%% concept/caution/example/takeaway boxes, \sectiondivider). Changes for this course:
%%   * course title block (\decktitle / \makedecktitle) instead of \makebeamertitle
%%   * \zh{...} typesets Chinese (book titles) under pdflatex via CJKutf8
%%   * researchbox for the "Research in the wild" frames
%%   * section pages off; TikZ libraries and the course map loaded here
%% Build with pdflatex only (two passes); tools/build_slides.py does this for you.

\documentclass[10pt,english,aspectratio=169]{beamer}

\usepackage[T1]{fontenc}
\usepackage{textcomp}
\usepackage[utf8]{inputenc}
\usepackage{babel}
\usepackage{CJKutf8}

\usepackage{amsmath, amssymb, amsthm, graphicx}
\usepackage{booktabs, multirow, tabularx}
\usepackage{tikz}
\usetikzlibrary{positioning,arrows.meta,shapes,calc,fit,backgrounds}
\usepackage{pgfplots}
\pgfplotsset{compat=1.16}
\usepackage[most]{tcolorbox}
\tcbuselibrary{skins, breakable}
\usepackage{listings}
\usepackage{xcolor}

\lstset{
  basicstyle=\ttfamily\small,
  keywordstyle=\color{blue!80!black}\bfseries,
  commentstyle=\color{green!50!black}\itshape,
  stringstyle=\color{orange!80!black},
  backgroundcolor=\color{gray!8},
  frame=single,
  rulecolor=\color{gray!40},
  breaklines=true,
  showstringspaces=false,
  numbers=none,
}

\usetheme{metropolis}
\usefonttheme{professionalfonts}
\metroset{sectionpage=none}

\definecolor{LightGray}{RGB}{230,230,230}
\definecolor{RedTitle}{RGB}{0,0,200}
\definecolor{VeryLightGray}{RGB}{245,245,245}
\definecolor{DarkText}{RGB}{0,0,0}
\definecolor{UCRBlue}{RGB}{0,45,106}
\definecolor{UCRGold}{RGB}{241,171,0}

% Stage colours of the course arc (used by the course map and elsewhere)
\colorlet{StageTrunk}{blue!12}
\colorlet{StageBranch}{cyan!14}
\colorlet{StageMachine}{gray!18}
\colorlet{StageWall}{red!14}
\colorlet{StageTools}{orange!20}
\colorlet{StageData}{green!16}
\colorlet{StageAgents}{violet!14}

\hypersetup{bookmarks=false,breaklinks=true,pdfborder={0 0 0},colorlinks=true,
  linkcolor=DarkText,urlcolor=RedTitle}

\setbeamercolor{background canvas}{bg=white}
\setbeamercolor{title}{fg=RedTitle}
\setbeamercolor{frametitle}{bg=VeryLightGray,fg=DarkText}
\setbeamercolor{normal text}{fg=DarkText}
\setbeamercolor{alerted text}{fg=RedTitle}
\setbeamersize{text margin left=5mm, text margin right=5mm}
\addtobeamertemplate{frametitle}{}{\vspace{-0.5em}}

\setbeamertemplate{navigation symbols}{}
\setbeamertemplate{footline}{%
  \hspace{5mm}{\usebeamerfont{page number in head/foot}\color{black!45}%
  ECON 282E \,$\cdot$\, \insertshortdeck}\hfill%
  \usebeamercolor[fg]{page number in head/foot}%
  \usebeamerfont{page number in head/foot}%
  \insertframenumber\,/\,\inserttotalframenumber\kern1em\vskip2pt%
}

% ---------------------------------------------------------------------------
% Chinese text under pdflatex (book titles etc.)
\newcommand{\zh}[1]{\begin{CJK*}{UTF8}{gbsn}#1\end{CJK*}}

% ---------------------------------------------------------------------------
% Deck title block
%   \decktitle{session}{deck id}{title}{date}
\newcommand{\insertshortdeck}{}
\newcommand{\decksession}{}
\newcommand{\deckid}{}
\newcommand{\decktitletext}{}
\newcommand{\deckdate}{}
\newcommand{\decktitle}[4]{%
  \renewcommand{\decksession}{#1}%
  \renewcommand{\deckid}{#2}%
  \renewcommand{\decktitletext}{#3}%
  \renewcommand{\deckdate}{#4}%
  \renewcommand{\insertshortdeck}{Session #1 \,$\cdot$\, Deck #1.#2}%
  \title{#3}%
  \author{Zhigang Feng}%
  \date{#4}%
}
\newcommand{\makedecktitle}{%
  \begin{frame}[plain,noframenumbering]
    \begin{tikzpicture}[remember picture,overlay]
      \fill[UCRBlue] (current page.north west) rectangle ([yshift=-0.55cm]current page.north east);
      \fill[UCRGold] ([yshift=-0.55cm]current page.north west) rectangle ([yshift=-0.62cm]current page.north east);
      \node[anchor=west, text=white, font=\small\bfseries] at ([xshift=5mm,yshift=-0.275cm]current page.north west)
        {ECON 282E \,$\cdot$\, Foundations of Macroeconomics};
      \node[anchor=east, text=white, font=\small] at ([xshift=-5mm,yshift=-0.275cm]current page.north east)
        {UC Riverside \,$\cdot$\, Fall 2026};
    \end{tikzpicture}
    \vspace*{1.1cm}
    {\usebeamerfont{title}\color{black!55}\large Session \decksession\ \,$\cdot$\, Deck \decksession.\deckid\par}
    \vspace{0.25cm}
    {\usebeamerfont{title}\color{RedTitle}\LARGE\bfseries \decktitletext\par}
    \vspace{0.7cm}
    {\large Zhigang Feng\par}
    \vspace{0.15cm}
    {\color{black!65}\deckdate\par}
    \vfill
    {\footnotesize\itshape\color{black!55}Build it on paper $\;\to\;$ solve it on the machine $\;\to\;$ push it to the frontier with AI\par}
  \end{frame}}

% ---------------------------------------------------------------------------
% Section divider (plain frame, not counted)
\newcommand{\sectiondivider}[2]{%
\begin{frame}[plain,noframenumbering]
\begin{center}
\vfill
{\Large\textbf{#1}\par}
\vspace{0.35cm}
{\normalsize #2\par}
\vfill
\end{center}
\end{frame}
}

% ---------------------------------------------------------------------------
% Boxes
\newtcolorbox{conceptbox}[1][]{colback=blue!3, colframe=RedTitle!55, boxrule=0.35mm,
  arc=1mm, left=2mm, right=2mm, top=1mm, bottom=1mm, #1}
\newtcolorbox{cautionbox}[1][]{colback=red!3, colframe=red!50!black, boxrule=0.35mm,
  arc=1mm, left=2mm, right=2mm, top=1mm, bottom=1mm, #1}
\newtcolorbox{examplebox}[1][]{colback=green!3, colframe=green!45!black, boxrule=0.35mm,
  arc=1mm, left=2mm, right=2mm, top=1mm, bottom=1mm, #1}
\newtcolorbox{takeawaybox}[1][]{colback=VeryLightGray, colframe=RedTitle!55, boxrule=0.35mm,
  arc=1mm, left=2mm, right=2mm, top=1mm, bottom=1mm, #1}
\newtcolorbox{researchbox}[1][]{enhanced, colback=violet!4, colframe=violet!60!black,
  boxrule=0.35mm, arc=1mm, left=2mm, right=2mm, top=1mm, bottom=1mm,
  fonttitle=\bfseries\small, title={Research in the wild}, #1}

\newcommand{\compacttables}{\renewcommand{\arraystretch}{1.15}}
\newcommand{\src}[1]{{\tiny\color{black!55}#1}}

% ---------------------------------------------------------------------------
\input{../preamble/course_map.tex}

%%   \decktitle{1}{A1}{Who I Am \& Why This Course}{September 24, 2026}
%%   \begin{document}
%%   \makedecktitle
%%   ...
%%
%% Adapted from Courses/AI-ECON-2026/source/shared_preamble.tex (Metropolis theme,
%% concept/caution/example/takeaway boxes, \sectiondivider). Changes for this course:
%%   * course title block (\decktitle / \makedecktitle) instead of \makebeamertitle
%%   * \zh{...} typesets Chinese (book titles) under pdflatex via CJKutf8
%%   * researchbox for the "Research in the wild" frames
%%   * section pages off; TikZ libraries and the course map loaded here
%% Build with pdflatex only (two passes); tools/build_slides.py does this for you.

\documentclass[10pt,english,aspectratio=169]{beamer}

\usepackage[T1]{fontenc}
\usepackage{textcomp}
\usepackage[utf8]{inputenc}
\usepackage{babel}
\usepackage{CJKutf8}

\usepackage{amsmath, amssymb, amsthm, graphicx}
\usepackage{booktabs, multirow, tabularx}
\usepackage{tikz}
\usetikzlibrary{positioning,arrows.meta,shapes,calc,fit,backgrounds}
\usepackage{pgfplots}
\pgfplotsset{compat=1.16}
\usepackage[most]{tcolorbox}
\tcbuselibrary{skins, breakable}
\usepackage{listings}
\usepackage{xcolor}

\lstset{
  basicstyle=\ttfamily\small,
  keywordstyle=\color{blue!80!black}\bfseries,
  commentstyle=\color{green!50!black}\itshape,
  stringstyle=\color{orange!80!black},
  backgroundcolor=\color{gray!8},
  frame=single,
  rulecolor=\color{gray!40},
  breaklines=true,
  showstringspaces=false,
  numbers=none,
}

\usetheme{metropolis}
\usefonttheme{professionalfonts}
\metroset{sectionpage=none}

\definecolor{LightGray}{RGB}{230,230,230}
\definecolor{RedTitle}{RGB}{0,0,200}
\definecolor{VeryLightGray}{RGB}{245,245,245}
\definecolor{DarkText}{RGB}{0,0,0}
\definecolor{UCRBlue}{RGB}{0,45,106}
\definecolor{UCRGold}{RGB}{241,171,0}

% Stage colours of the course arc (used by the course map and elsewhere)
\colorlet{StageTrunk}{blue!12}
\colorlet{StageBranch}{cyan!14}
\colorlet{StageMachine}{gray!18}
\colorlet{StageWall}{red!14}
\colorlet{StageTools}{orange!20}
\colorlet{StageData}{green!16}
\colorlet{StageAgents}{violet!14}

\hypersetup{bookmarks=false,breaklinks=true,pdfborder={0 0 0},colorlinks=true,
  linkcolor=DarkText,urlcolor=RedTitle}

\setbeamercolor{background canvas}{bg=white}
\setbeamercolor{title}{fg=RedTitle}
\setbeamercolor{frametitle}{bg=VeryLightGray,fg=DarkText}
\setbeamercolor{normal text}{fg=DarkText}
\setbeamercolor{alerted text}{fg=RedTitle}
\setbeamersize{text margin left=5mm, text margin right=5mm}
\addtobeamertemplate{frametitle}{}{\vspace{-0.5em}}

\setbeamertemplate{navigation symbols}{}
\setbeamertemplate{footline}{%
  \hspace{5mm}{\usebeamerfont{page number in head/foot}\color{black!45}%
  ECON 282E \,$\cdot$\, \insertshortdeck}\hfill%
  \usebeamercolor[fg]{page number in head/foot}%
  \usebeamerfont{page number in head/foot}%
  \insertframenumber\,/\,\inserttotalframenumber\kern1em\vskip2pt%
}

% ---------------------------------------------------------------------------
% Chinese text under pdflatex (book titles etc.)
\newcommand{\zh}[1]{\begin{CJK*}{UTF8}{gbsn}#1\end{CJK*}}

% ---------------------------------------------------------------------------
% Deck title block
%   \decktitle{session}{deck id}{title}{date}
\newcommand{\insertshortdeck}{}
\newcommand{\decksession}{}
\newcommand{\deckid}{}
\newcommand{\decktitletext}{}
\newcommand{\deckdate}{}
\newcommand{\decktitle}[4]{%
  \renewcommand{\decksession}{#1}%
  \renewcommand{\deckid}{#2}%
  \renewcommand{\decktitletext}{#3}%
  \renewcommand{\deckdate}{#4}%
  \renewcommand{\insertshortdeck}{Session #1 \,$\cdot$\, Deck #1.#2}%
  \title{#3}%
  \author{Zhigang Feng}%
  \date{#4}%
}
\newcommand{\makedecktitle}{%
  \begin{frame}[plain,noframenumbering]
    \begin{tikzpicture}[remember picture,overlay]
      \fill[UCRBlue] (current page.north west) rectangle ([yshift=-0.55cm]current page.north east);
      \fill[UCRGold] ([yshift=-0.55cm]current page.north west) rectangle ([yshift=-0.62cm]current page.north east);
      \node[anchor=west, text=white, font=\small\bfseries] at ([xshift=5mm,yshift=-0.275cm]current page.north west)
        {ECON 282E \,$\cdot$\, Foundations of Macroeconomics};
      \node[anchor=east, text=white, font=\small] at ([xshift=-5mm,yshift=-0.275cm]current page.north east)
        {UC Riverside \,$\cdot$\, Fall 2026};
    \end{tikzpicture}
    \vspace*{1.1cm}
    {\usebeamerfont{title}\color{black!55}\large Session \decksession\ \,$\cdot$\, Deck \decksession.\deckid\par}
    \vspace{0.25cm}
    {\usebeamerfont{title}\color{RedTitle}\LARGE\bfseries \decktitletext\par}
    \vspace{0.7cm}
    {\large Zhigang Feng\par}
    \vspace{0.15cm}
    {\color{black!65}\deckdate\par}
    \vfill
    {\footnotesize\itshape\color{black!55}Build it on paper $\;\to\;$ solve it on the machine $\;\to\;$ push it to the frontier with AI\par}
  \end{frame}}

% ---------------------------------------------------------------------------
% Section divider (plain frame, not counted)
\newcommand{\sectiondivider}[2]{%
\begin{frame}[plain,noframenumbering]
\begin{center}
\vfill
{\Large\textbf{#1}\par}
\vspace{0.35cm}
{\normalsize #2\par}
\vfill
\end{center}
\end{frame}
}

% ---------------------------------------------------------------------------
% Boxes
\newtcolorbox{conceptbox}[1][]{colback=blue!3, colframe=RedTitle!55, boxrule=0.35mm,
  arc=1mm, left=2mm, right=2mm, top=1mm, bottom=1mm, #1}
\newtcolorbox{cautionbox}[1][]{colback=red!3, colframe=red!50!black, boxrule=0.35mm,
  arc=1mm, left=2mm, right=2mm, top=1mm, bottom=1mm, #1}
\newtcolorbox{examplebox}[1][]{colback=green!3, colframe=green!45!black, boxrule=0.35mm,
  arc=1mm, left=2mm, right=2mm, top=1mm, bottom=1mm, #1}
\newtcolorbox{takeawaybox}[1][]{colback=VeryLightGray, colframe=RedTitle!55, boxrule=0.35mm,
  arc=1mm, left=2mm, right=2mm, top=1mm, bottom=1mm, #1}
\newtcolorbox{researchbox}[1][]{enhanced, colback=violet!4, colframe=violet!60!black,
  boxrule=0.35mm, arc=1mm, left=2mm, right=2mm, top=1mm, bottom=1mm,
  fonttitle=\bfseries\small, title={Research in the wild}, #1}

\newcommand{\compacttables}{\renewcommand{\arraystretch}{1.15}}
\newcommand{\src}[1]{{\tiny\color{black!55}#1}}

% ---------------------------------------------------------------------------
%% Course map: the recurring "you are here" slide for ECON 282E.
%% Loaded by course_preamble.tex. Pattern follows AI-ECON-2026 lec01_map.tex, but the map
%% shows the whole ten-session course rather than one lecture.
%%
%%   \CourseMap{s}{label}   s = 1..10 highlights session s and prints "you are here: label"
%%                          (e.g. \CourseMap{1}{1.A2}); earlier sessions get a check mark.
%%   \CourseMap{0}{}        full overview, nothing highlighted.
%%
%% Note: TikZ option lists are comma-split before TeX conditionals run, so every \ifnum
%% below sits outside [...] option lists.

\newcommand{\CourseMap}[2]{%
\begin{center}
\begin{tikzpicture}[
    sess/.style={rectangle, rounded corners=2pt, draw=black!45, minimum width=1.34cm,
        minimum height=1.12cm, text width=1.26cm, align=center, inner sep=1pt},
    stage/.style={font=\tiny\bfseries, text=black!70},
]
  \foreach \i/\d/\t/\c in {%
      1/Sep 24/Intro \\ Growth/StageTrunk,
      2/Oct 1/HA $\cdot$ OLG \\ Policy/StageBranch,
      3/Oct 8/Num.\ DP \\ Python/StageMachine,
      4/Oct 15/Numerics \\ I/StageMachine,
      5/Oct 22/Numerics \\ II/StageMachine,
      6/Oct 29/The wall \\ \textbf{Midterm}/StageWall,
      7/Nov 5/ML \\ Deep nets/StageTools,
      8/Nov 12/RL \\ HA $+$ DL/StageTools,
      9/Nov 19/Text \\ LLMs/StageData,
      10/Dec 3/Agents \\ \textbf{Projects}/StageAgents} {
    \node[sess, fill=\c] (n\i) at ({1.46*(\i-1)}, 0)
      {{\scriptsize\bfseries S\i}\\[-1pt]{\tiny\color{black!60}\d}\\[1pt]{\tiny \t}};
    \ifnum\i<#1
      \node[font=\scriptsize, text=green!45!black] at ([xshift=-2.5mm,yshift=-2.2mm]n\i.north east) {$\checkmark$};
    \fi
    \ifnum\i=#1
      \draw[RedTitle, very thick, rounded corners=2pt]
        ([xshift=-0.6mm,yshift=-0.6mm]n\i.south west) rectangle ([xshift=0.6mm,yshift=0.6mm]n\i.north east);
      % keep the label inside the picture at both ends of the row
      \ifnum\i<3
        \node[font=\scriptsize\bfseries, text=RedTitle, anchor=south west] at ([yshift=1.2mm]n\i.north west)
          {$\blacktriangledown$\,you are here: #2};
      \else\ifnum\i>8
        \node[font=\scriptsize\bfseries, text=RedTitle, anchor=south east] at ([yshift=1.2mm]n\i.north east)
          {you are here: #2\,$\blacktriangledown$};
      \else
        \node[font=\scriptsize\bfseries, text=RedTitle, anchor=south] at ([yshift=1.2mm]n\i.north)
          {you are here: #2\,$\blacktriangledown$};
      \fi\fi
    \fi
  }
  % stage band under the sessions
  \foreach \a/\b/\lab/\c in {1/1/trunk/StageTrunk, 2/2/branches/StageBranch,
      3/5/the machine $\cdot$ classical methods/StageMachine, 6/6/the wall/StageWall,
      7/8/new tools/StageTools, 9/9/new data/StageData, 10/10/colleagues/StageAgents} {
    \fill[\c] ([yshift=-1.5mm]n\a.south west) rectangle ([yshift=-4.6mm]n\b.south east);
    \path ([yshift=-1.5mm]n\a.south west) -- ([yshift=-4.6mm]n\b.south east)
      node[midway, stage] {\lab};
  }
  \node[font=\footnotesize\itshape, text=black!70] at ([yshift=-9.5mm]$(n1.south)!0.5!(n10.south)$)
    {Build it on paper $\;\to\;$ solve it on the machine $\;\to\;$ push it to the frontier with AI};
\end{tikzpicture}
\end{center}
}


%%   \decktitle{1}{A1}{Who I Am \& Why This Course}{September 24, 2026}
%%   \begin{document}
%%   \makedecktitle
%%   ...
%%
%% Adapted from Courses/AI-ECON-2026/source/shared_preamble.tex (Metropolis theme,
%% concept/caution/example/takeaway boxes, \sectiondivider). Changes for this course:
%%   * course title block (\decktitle / \makedecktitle) instead of \makebeamertitle
%%   * \zh{...} typesets Chinese (book titles) under pdflatex via CJKutf8
%%   * researchbox for the "Research in the wild" frames
%%   * section pages off; TikZ libraries and the course map loaded here
%% Build with pdflatex only (two passes); tools/build_slides.py does this for you.

\documentclass[10pt,english,aspectratio=169]{beamer}

\usepackage[T1]{fontenc}
\usepackage{textcomp}
\usepackage[utf8]{inputenc}
\usepackage{babel}
\usepackage{CJKutf8}

\usepackage{amsmath, amssymb, amsthm, graphicx}
\usepackage{booktabs, multirow, tabularx}
\usepackage{tikz}
\usetikzlibrary{positioning,arrows.meta,shapes,calc,fit,backgrounds}
\usepackage{pgfplots}
\pgfplotsset{compat=1.16}
\usepackage[most]{tcolorbox}
\tcbuselibrary{skins, breakable}
\usepackage{listings}
\usepackage{xcolor}

\lstset{
  basicstyle=\ttfamily\small,
  keywordstyle=\color{blue!80!black}\bfseries,
  commentstyle=\color{green!50!black}\itshape,
  stringstyle=\color{orange!80!black},
  backgroundcolor=\color{gray!8},
  frame=single,
  rulecolor=\color{gray!40},
  breaklines=true,
  showstringspaces=false,
  numbers=none,
}

\usetheme{metropolis}
\usefonttheme{professionalfonts}
\metroset{sectionpage=none}

\definecolor{LightGray}{RGB}{230,230,230}
\definecolor{RedTitle}{RGB}{0,0,200}
\definecolor{VeryLightGray}{RGB}{245,245,245}
\definecolor{DarkText}{RGB}{0,0,0}
\definecolor{UCRBlue}{RGB}{0,45,106}
\definecolor{UCRGold}{RGB}{241,171,0}

% Stage colours of the course arc (used by the course map and elsewhere)
\colorlet{StageTrunk}{blue!12}
\colorlet{StageBranch}{cyan!14}
\colorlet{StageMachine}{gray!18}
\colorlet{StageWall}{red!14}
\colorlet{StageTools}{orange!20}
\colorlet{StageData}{green!16}
\colorlet{StageAgents}{violet!14}

\hypersetup{bookmarks=false,breaklinks=true,pdfborder={0 0 0},colorlinks=true,
  linkcolor=DarkText,urlcolor=RedTitle}

\setbeamercolor{background canvas}{bg=white}
\setbeamercolor{title}{fg=RedTitle}
\setbeamercolor{frametitle}{bg=VeryLightGray,fg=DarkText}
\setbeamercolor{normal text}{fg=DarkText}
\setbeamercolor{alerted text}{fg=RedTitle}
\setbeamersize{text margin left=5mm, text margin right=5mm}
\addtobeamertemplate{frametitle}{}{\vspace{-0.5em}}

\setbeamertemplate{navigation symbols}{}
\setbeamertemplate{footline}{%
  \hspace{5mm}{\usebeamerfont{page number in head/foot}\color{black!45}%
  ECON 282E \,$\cdot$\, \insertshortdeck}\hfill%
  \usebeamercolor[fg]{page number in head/foot}%
  \usebeamerfont{page number in head/foot}%
  \insertframenumber\,/\,\inserttotalframenumber\kern1em\vskip2pt%
}

% ---------------------------------------------------------------------------
% Chinese text under pdflatex (book titles etc.)
\newcommand{\zh}[1]{\begin{CJK*}{UTF8}{gbsn}#1\end{CJK*}}

% ---------------------------------------------------------------------------
% Deck title block
%   \decktitle{session}{deck id}{title}{date}
\newcommand{\insertshortdeck}{}
\newcommand{\decksession}{}
\newcommand{\deckid}{}
\newcommand{\decktitletext}{}
\newcommand{\deckdate}{}
\newcommand{\decktitle}[4]{%
  \renewcommand{\decksession}{#1}%
  \renewcommand{\deckid}{#2}%
  \renewcommand{\decktitletext}{#3}%
  \renewcommand{\deckdate}{#4}%
  \renewcommand{\insertshortdeck}{Session #1 \,$\cdot$\, Deck #1.#2}%
  \title{#3}%
  \author{Zhigang Feng}%
  \date{#4}%
}
\newcommand{\makedecktitle}{%
  \begin{frame}[plain,noframenumbering]
    \begin{tikzpicture}[remember picture,overlay]
      \fill[UCRBlue] (current page.north west) rectangle ([yshift=-0.55cm]current page.north east);
      \fill[UCRGold] ([yshift=-0.55cm]current page.north west) rectangle ([yshift=-0.62cm]current page.north east);
      \node[anchor=west, text=white, font=\small\bfseries] at ([xshift=5mm,yshift=-0.275cm]current page.north west)
        {ECON 282E \,$\cdot$\, Foundations of Macroeconomics};
      \node[anchor=east, text=white, font=\small] at ([xshift=-5mm,yshift=-0.275cm]current page.north east)
        {UC Riverside \,$\cdot$\, Fall 2026};
    \end{tikzpicture}
    \vspace*{1.1cm}
    {\usebeamerfont{title}\color{black!55}\large Session \decksession\ \,$\cdot$\, Deck \decksession.\deckid\par}
    \vspace{0.25cm}
    {\usebeamerfont{title}\color{RedTitle}\LARGE\bfseries \decktitletext\par}
    \vspace{0.7cm}
    {\large Zhigang Feng\par}
    \vspace{0.15cm}
    {\color{black!65}\deckdate\par}
    \vfill
    {\footnotesize\itshape\color{black!55}Build it on paper $\;\to\;$ solve it on the machine $\;\to\;$ push it to the frontier with AI\par}
  \end{frame}}

% ---------------------------------------------------------------------------
% Section divider (plain frame, not counted)
\newcommand{\sectiondivider}[2]{%
\begin{frame}[plain,noframenumbering]
\begin{center}
\vfill
{\Large\textbf{#1}\par}
\vspace{0.35cm}
{\normalsize #2\par}
\vfill
\end{center}
\end{frame}
}

% ---------------------------------------------------------------------------
% Boxes
\newtcolorbox{conceptbox}[1][]{colback=blue!3, colframe=RedTitle!55, boxrule=0.35mm,
  arc=1mm, left=2mm, right=2mm, top=1mm, bottom=1mm, #1}
\newtcolorbox{cautionbox}[1][]{colback=red!3, colframe=red!50!black, boxrule=0.35mm,
  arc=1mm, left=2mm, right=2mm, top=1mm, bottom=1mm, #1}
\newtcolorbox{examplebox}[1][]{colback=green!3, colframe=green!45!black, boxrule=0.35mm,
  arc=1mm, left=2mm, right=2mm, top=1mm, bottom=1mm, #1}
\newtcolorbox{takeawaybox}[1][]{colback=VeryLightGray, colframe=RedTitle!55, boxrule=0.35mm,
  arc=1mm, left=2mm, right=2mm, top=1mm, bottom=1mm, #1}
\newtcolorbox{researchbox}[1][]{enhanced, colback=violet!4, colframe=violet!60!black,
  boxrule=0.35mm, arc=1mm, left=2mm, right=2mm, top=1mm, bottom=1mm,
  fonttitle=\bfseries\small, title={Research in the wild}, #1}

\newcommand{\compacttables}{\renewcommand{\arraystretch}{1.15}}
\newcommand{\src}[1]{{\tiny\color{black!55}#1}}

% ---------------------------------------------------------------------------
%% Course map: the recurring "you are here" slide for ECON 282E.
%% Loaded by course_preamble.tex. Pattern follows AI-ECON-2026 lec01_map.tex, but the map
%% shows the whole ten-session course rather than one lecture.
%%
%%   \CourseMap{s}{label}   s = 1..10 highlights session s and prints "you are here: label"
%%                          (e.g. \CourseMap{1}{1.A2}); earlier sessions get a check mark.
%%   \CourseMap{0}{}        full overview, nothing highlighted.
%%
%% Note: TikZ option lists are comma-split before TeX conditionals run, so every \ifnum
%% below sits outside [...] option lists.

\newcommand{\CourseMap}[2]{%
\begin{center}
\begin{tikzpicture}[
    sess/.style={rectangle, rounded corners=2pt, draw=black!45, minimum width=1.34cm,
        minimum height=1.12cm, text width=1.26cm, align=center, inner sep=1pt},
    stage/.style={font=\tiny\bfseries, text=black!70},
]
  \foreach \i/\d/\t/\c in {%
      1/Sep 24/Intro \\ Growth/StageTrunk,
      2/Oct 1/HA $\cdot$ OLG \\ Policy/StageBranch,
      3/Oct 8/Num.\ DP \\ Python/StageMachine,
      4/Oct 15/Numerics \\ I/StageMachine,
      5/Oct 22/Numerics \\ II/StageMachine,
      6/Oct 29/The wall \\ \textbf{Midterm}/StageWall,
      7/Nov 5/ML \\ Deep nets/StageTools,
      8/Nov 12/RL \\ HA $+$ DL/StageTools,
      9/Nov 19/Text \\ LLMs/StageData,
      10/Dec 3/Agents \\ \textbf{Projects}/StageAgents} {
    \node[sess, fill=\c] (n\i) at ({1.46*(\i-1)}, 0)
      {{\scriptsize\bfseries S\i}\\[-1pt]{\tiny\color{black!60}\d}\\[1pt]{\tiny \t}};
    \ifnum\i<#1
      \node[font=\scriptsize, text=green!45!black] at ([xshift=-2.5mm,yshift=-2.2mm]n\i.north east) {$\checkmark$};
    \fi
    \ifnum\i=#1
      \draw[RedTitle, very thick, rounded corners=2pt]
        ([xshift=-0.6mm,yshift=-0.6mm]n\i.south west) rectangle ([xshift=0.6mm,yshift=0.6mm]n\i.north east);
      % keep the label inside the picture at both ends of the row
      \ifnum\i<3
        \node[font=\scriptsize\bfseries, text=RedTitle, anchor=south west] at ([yshift=1.2mm]n\i.north west)
          {$\blacktriangledown$\,you are here: #2};
      \else\ifnum\i>8
        \node[font=\scriptsize\bfseries, text=RedTitle, anchor=south east] at ([yshift=1.2mm]n\i.north east)
          {you are here: #2\,$\blacktriangledown$};
      \else
        \node[font=\scriptsize\bfseries, text=RedTitle, anchor=south] at ([yshift=1.2mm]n\i.north)
          {you are here: #2\,$\blacktriangledown$};
      \fi\fi
    \fi
  }
  % stage band under the sessions
  \foreach \a/\b/\lab/\c in {1/1/trunk/StageTrunk, 2/2/branches/StageBranch,
      3/5/the machine $\cdot$ classical methods/StageMachine, 6/6/the wall/StageWall,
      7/8/new tools/StageTools, 9/9/new data/StageData, 10/10/colleagues/StageAgents} {
    \fill[\c] ([yshift=-1.5mm]n\a.south west) rectangle ([yshift=-4.6mm]n\b.south east);
    \path ([yshift=-1.5mm]n\a.south west) -- ([yshift=-4.6mm]n\b.south east)
      node[midway, stage] {\lab};
  }
  \node[font=\footnotesize\itshape, text=black!70] at ([yshift=-9.5mm]$(n1.south)!0.5!(n10.south)$)
    {Build it on paper $\;\to\;$ solve it on the machine $\;\to\;$ push it to the frontier with AI};
\end{tikzpicture}
\end{center}
}


\graphicspath{{./figures/}}

\lstset{language=Python, basicstyle=\ttfamily\scriptsize, frame=none,
        backgroundcolor=\color{gray!8}, aboveskip=2pt, belowskip=2pt}

\decktitle{3}{A4}{Worked Example: The RBC Model, Fully Discretized}{Thursday, October 8, 2026}

\begin{document}

\makedecktitle

%=============================================================================
\begin{frame}{Where We Are}
\CourseMap{3}{3.A4}
\vspace{-1mm}
\begin{center}
\small \textbf{Block A:} \textbf{3.A1} what it means to solve a model $\;\cdot\;$ \textbf{3.A2} value function iteration on a grid $\;\cdot\;$ \textbf{3.A3} Euler equations and time iteration $\;\cdot\;$ \textbf{3.A4} the worked example, end to end.
\end{center}
\end{frame}

%=============================================================================
\begin{frame}{The Model, and the One Thing We Know About It Exactly}
\begin{columns}[T]
\begin{column}{0.52\textwidth}
\small
The standard deterministic RBC problem, now with \textbf{realistic depreciation}:
\[
  \max \sum_{t=0}^{\infty}\beta^{t}\ln c_t
  \quad\text{s.t.}\quad c_t+k_{t+1}=Ak_t^{\alpha}+(1-\delta)k_t .
\]
Its steady state follows from $1=\beta\big(\alpha Ak^{\alpha-1}+1-\delta\big)$:
\[
  k^{*}=\Big(\frac{\alpha\beta A}{1-\beta(1-\delta)}\Big)^{\frac{1}{1-\alpha}} .
\]
At $A=1$, $\alpha=0.36$, $\beta=0.95$, $\delta=0.10$: \;$k^{*}=\mathbf{3.8219}$, $c^{*}=1.2382$.
\end{column}
\begin{column}{0.45\textwidth}
\begin{cautionbox}[title=\textbf{What changed from this morning}]\small
$\delta=0.10$, not 1. So there is a closed-form \textbf{steady state} --- but \emph{no} closed-form \textbf{policy}.
\end{cautionbox}
\medskip
\begin{takeawaybox}\footnotesize
That is the ordinary situation, and it is why this example is worth doing. You will almost never have $g(k)$ to check against. You will very often have \textbf{one point} on it, and one point is enough to catch most mistakes.
\end{takeawaybox}
\end{column}
\end{columns}
\vspace{1mm}
\src{Version 0 of the RBC ladder in QM Lec.~3A; the $\delta=1$ special case of 3.A1--3.A3 is the one with a closed-form policy.}
\end{frame}

%=============================================================================
\begin{frame}{``Fully Discretized'' --- What That Commits Us To}
\begin{columns}[T]
\begin{column}{0.53\textwidth}
\small
One decision does all the work: \textbf{the choice set is the state grid}.
\[
  k'\in\{k_1,\dots,k_N\}=\text{the same grid } k \text{ lives on.}
\]
\begin{itemize}\footnotesize
  \item No \textbf{interpolation}: $v$ is only ever evaluated \emph{at} a grid point, because $k'$ is always a grid point.
  \item No \textbf{maximization routine}: the $\max$ is a search over $N$ numbers, so \texttt{argmax} does it.
  \item No \textbf{root-finder}, no derivatives, no Jacobian.
\end{itemize}
\end{column}
\begin{column}{0.44\textwidth}
\begin{conceptbox}[title=\textbf{Why start here}]\small
It is the only version of the algorithm with \textbf{no numerical method inside it}. Everything is indexing and arithmetic --- so if the answer is wrong, the economics or the code is wrong, and nothing else can be blamed.
\end{conceptbox}
\medskip
\begin{cautionbox}\footnotesize
The price is the staircase of 3.A2, and an $O(N^2)$ inner loop. Next week buys both back.
\end{cautionbox}
\end{column}
\end{columns}
\end{frame}

%=============================================================================
\begin{frame}[fragile]{Step 1: Parameters, the Benchmark, and the Grid}
\begin{columns}[T]
\begin{column}{0.55\textwidth}
\begin{lstlisting}
import numpy as np

A, alpha, beta, delta = 1.0, 0.36, 0.95, 0.10

# the one thing we know exactly
kss = (alpha*beta*A /
       (1 - beta*(1-delta)))**(1/(1-alpha))
# 3.8218909152179084

# the grid IS the choice set
N  = 500
kg = np.linspace(0.25*kss, 1.75*kss, N)
h  = kg[1] - kg[0]          # 0.01149
\end{lstlisting}
\end{column}
\begin{column}{0.42\textwidth}
\small
\begin{itemize}\footnotesize
  \item The benchmark is computed \textbf{first}, before any solver exists. If you write it afterwards you will unconsciously write it to agree.
  \item The grid is centred on $k^{*}$ and runs $\pm75\%$ --- wide enough that the policy never touches an edge.
  \item $N=500$ is a choice, not a fact. We will make it earn its keep.
\end{itemize}
\end{column}
\end{columns}
\end{frame}

%=============================================================================
\begin{frame}[fragile]{Step 2: The Return Matrix, Built Once}
\begin{columns}[T]
\begin{column}{0.55\textwidth}
\begin{lstlisting}
# cash on hand at each state: (N,)
y = A*kg**alpha + (1-delta)*kg

# consumption for every (state, choice): (N,N)
C = y[:, None] - kg[None, :]

# u(c), with -1e10 where c <= 0
U = np.full_like(C, -1e10)
np.log(C, out=U, where=C > 0)
\end{lstlisting}
\vspace{-1mm}
\footnotesize
\texttt{U[i, j]} $=u\big(Ak_i^{\alpha}+(1-\delta)k_i-k_j\big)$.
\end{column}
\begin{column}{0.42\textwidth}
\begin{conceptbox}[title=\textbf{Two ideas in six lines}]\small
\textbf{Broadcasting} builds every state against every choice at once --- 3.B2's whole subject.\\[1mm]
\textbf{$U$ does not contain $v$}, so it is built once and reused every sweep.
\end{conceptbox}
\medskip
\begin{cautionbox}\footnotesize
The feasibility penalty must be a \emph{number}, not $-\infty$: $-\infty+\beta v$ is $-\infty$, but $-\infty-\infty$ in a later comparison can produce \texttt{nan}, and \texttt{nan} loses every comparison silently.
\end{cautionbox}
\end{column}
\end{columns}
\end{frame}

%=============================================================================
\begin{frame}[fragile]{Step 3: One Bellman Step Is Three Lines}
\begin{columns}[T]
\begin{column}{0.54\textwidth}
\begin{lstlisting}
M   = U + beta*V[None, :]   # (N,N)
Vn  = M.max(axis=1)         # best value
pol = M.argmax(axis=1)      # best index
\end{lstlisting}
\vspace{-1mm}
\small
Read the shapes and the equation is there:
\[
  \underbrace{M_{ij}}_{(N,N)}=\underbrace{U_{ij}}_{\text{payoff}}+\beta\underbrace{v_j}_{\text{continuation}},
  \qquad v^{\text{new}}_i=\max_j M_{ij}.
\]
The maximisation is a \texttt{max} along an axis. That is the entire numerical content of this method.
\end{column}
\begin{column}{0.43\textwidth}
\begin{cautionbox}[title=\textbf{The axis is the whole game}]\footnotesize
\texttt{axis=1} reduces over \emph{choices} and leaves one value per state. \texttt{axis=0} would reduce over states and return something with the right shape and no meaning.\\[1mm]
Broadcasting will not warn you. \textbf{Assert the shape.}
\end{cautionbox}
\medskip
\begin{examplebox}\footnotesize
\texttt{argmax} returns an \emph{index}. Keeping the index rather than the value is what lets us read the policy off at the end.
\end{examplebox}
\end{column}
\end{columns}
\end{frame}

%=============================================================================
\begin{frame}[fragile]{Step 4: Iterate to the Fixed Point}
\begin{columns}[T]
\begin{column}{0.53\textwidth}
\begin{lstlisting}
V, it = np.zeros(N), 0
while True:
    M = U + beta*V[None, :]
    Vn  = M.max(axis=1)
    pol = M.argmax(axis=1)
    d = np.max(np.abs(Vn - V))
    V, it = Vn, it + 1
    if d < 1e-6:
        break

print(it)          # 228
\end{lstlisting}
\end{column}
\begin{column}{0.44\textwidth}
\small
\begin{center}\footnotesize
\begin{tabular}{@{}lr@{}}
\toprule
sweeps & 228 \\
$\ln\varepsilon/\ln\beta$ & 269 \\
wall time & 0.05\,s \\
comparisons & 57 million \\
\bottomrule
\end{tabular}
\end{center}
\medskip
\begin{examplebox}\footnotesize
\textbf{228, not 269.} The bound of 3.A2 assumed the worst case $\lVert v_1-v_0\rVert=1$; starting from $v_0=0$ the first step is already informative, so it finishes early. The bound is an upper bound, and it is allowed to be loose.
\end{examplebox}
\end{column}
\end{columns}
\end{frame}

%=============================================================================
\begin{frame}[fragile]{Step 5: Read the Policy Off the Indices}
\begin{columns}[T]
\begin{column}{0.51\textwidth}
\begin{lstlisting}
g = kg[pol]      # k'(k), on the grid
c = y - g        # c(k), by the budget

# the check we can actually run
i = np.argmin(np.abs(g - kg))
print(kg[i], kss)
# 3.8161465901950360 3.8218909152179084
\end{lstlisting}
\vspace{-1mm}
\small
\texttt{kg[pol]} is \textbf{fancy indexing}: \texttt{pol} is an array of indices, so this looks up all $N$ of them at once.
\end{column}
\begin{column}{0.46\textwidth}
\begin{conceptbox}[title=\textbf{Consumption comes from the budget}]\small
Never solve for $c$ separately. Recovering it from $c=y-g$ makes the budget constraint hold \textbf{exactly}, by construction, rather than to within a tolerance.
\end{conceptbox}
\medskip
\begin{takeawaybox}\footnotesize
This is the whole solver: \textbf{five steps, about twenty-five lines}, no numerical library beyond NumPy, and not one routine we have not explained.
\end{takeawaybox}
\end{column}
\end{columns}
\end{frame}

%=============================================================================
\begin{frame}{Does It Work? The 45-Degree Test}
\begin{columns}[c]
\begin{column}{0.52\textwidth}
\centering
\begin{tikzpicture}
\begin{axis}[
    width=7.7cm, height=5.4cm, xmin=0.9, xmax=6.8, ymin=0.9, ymax=6.8,
    xlabel={\scriptsize capital $k$}, ylabel={\scriptsize policy $k'$},
    tick label style={font=\scriptsize},
    axis x line*=bottom, axis y line*=left,
    legend style={font=\tiny, draw=none, fill=none, at={(0.02,0.97)}, anchor=north west}]
  \addplot[very thick, RedTitle] coordinates {(0.9555,1.2542) (1.0704,1.3691) (1.1852,1.4725) (1.3001,1.5873) (1.4150,1.6907) (1.5299,1.7941) (1.6448,1.9090) (1.7597,2.0124) (1.8746,2.1158) (1.9895,2.2192) (2.1043,2.3226) (2.2192,2.4145) (2.3341,2.5179) (2.4490,2.6213) (2.5639,2.7247) (2.6788,2.8281) (2.7937,2.9315) (2.9085,3.0234) (3.0234,3.1268) (3.1383,3.2302) (3.2532,3.3336) (3.3681,3.4255) (3.4830,3.5289) (3.5979,3.6323) (3.7127,3.7242) (3.8276,3.8276) (3.9425,3.9310) (4.0574,4.0229) (4.1723,4.1263) (4.2872,4.2182) (4.4021,4.3216) (4.5170,4.4136) (4.6318,4.5170) (4.7467,4.6204) (4.8616,4.7123) (4.9765,4.8157) (5.0914,4.9076) (5.2063,5.0110) (5.3212,5.1029) (5.4360,5.2063) (5.5509,5.2982) (5.6658,5.4016) (5.7807,5.4935) (5.8956,5.5969) (6.0105,5.6888) (6.1254,5.7922) (6.2403,5.8841) (6.3551,5.9875) (6.4700,6.0794) (6.5849,6.1828)};
  \addlegendentry{$g(k)$, computed}
  \addplot[thick, black!45, dashed] coordinates {(0.9555,0.9555) (6.6883,6.6883)};
  \addlegendentry{45$^\circ$}
  \addplot[only marks, mark=*, mark size=1.6pt, orange!85!black] coordinates {(3.8219,3.8219)};
  \addlegendentry{$k^{*}$, closed form}
\end{axis}
\end{tikzpicture}
\end{column}
\begin{column}{0.45\textwidth}
\small
A steady state is a fixed point of the policy: $g(k^{*})=k^{*}$. So the policy must cross the 45-degree line, and it must cross it \textbf{at the number we computed in Step 1}.
\medskip
\begin{center}\footnotesize
\begin{tabular}{@{}lr@{}}
\toprule
$k^{*}$, closed form & 3.8219 \\
crossing, computed & 3.8161 \\
difference & $-0.0057$ \\
grid spacing $h$ & 0.0115 \\
\bottomrule
\end{tabular}
\end{center}
\vspace{1mm}
\begin{examplebox}\footnotesize
The crossing is the \textbf{nearest grid point to $k^{*}$} --- it cannot be anything else. The solver is as accurate as the grid permits and no more, which is exactly the right answer.
\end{examplebox}
\end{column}
\end{columns}
\end{frame}

%=============================================================================
\begin{frame}{The Honest Verdict: Euler Residuals}
\begin{columns}[c]
\begin{column}{0.52\textwidth}
\centering
\begin{tikzpicture}
\begin{axis}[
    width=7.7cm, height=5.4cm, xmin=0.9, xmax=6.8, ymin=-4.2, ymax=-1.4,
    xlabel={\scriptsize capital $k$ (off the grid)}, ylabel={\scriptsize $\log_{10}\mathcal{E}$},
    tick label style={font=\scriptsize},
    axis x line*=bottom, axis y line*=left]
  \addplot[thin, RedTitle] coordinates {(0.9785,-2.0337) (1.0251,-2.0177) (1.0717,-2.1784) (1.1183,-2.1872) (1.1649,-2.1847) (1.2115,-1.8355) (1.2581,-2.8446) (1.3047,-2.6850) (1.3514,-2.5579) (1.3980,-2.8884) (1.4446,-3.5138) (1.4912,-3.1518) (1.5378,-2.6066) (1.5844,-2.8920) (1.6310,-1.9564) (1.6777,-2.5825) (1.7243,-2.8834) (1.7709,-1.8730) (1.8175,-2.6899) (1.8641,-3.2978) (1.9107,-2.6222) (1.9573,-3.0124) (2.0040,-2.0457) (2.0506,-3.0123) (2.0972,-2.2902) (2.1438,-3.1364) (2.1904,-3.1543) (2.2370,-3.5370) (2.2836,-2.9410) (2.3302,-1.9432) (2.3769,-2.7580) (2.4235,-2.1272) (2.4701,-2.6065) (2.5167,-2.4124) (2.5633,-2.5598) (2.6099,-2.3287) (2.6565,-2.2142) (2.7032,-2.2519) (2.7498,-2.2475) (2.7964,-2.1820) (2.8430,-2.3536) (2.8896,-2.1186) (2.9362,-2.4998) (2.9828,-2.0608) (3.0295,-2.7292) (3.0761,-2.0081) (3.1227,-3.2658) (3.1693,-2.1435) (3.2159,-3.0970) (3.2625,-2.1548) (3.3091,-2.6664) (3.3558,-2.1698) (3.4024,-2.4532) (3.4490,-2.2187) (3.4956,-2.3103) (3.5422,-2.3444) (3.5888,-2.2027) (3.6354,-2.5215) (3.6820,-2.1166) (3.7287,-2.8237) (3.7753,-3.4495) (3.8219,-7.5503) (3.8685,-3.4278) (3.9151,-2.8264) (3.9617,-2.1423) (4.0083,-2.5268) (4.0550,-2.2505) (4.1016,-2.3523) (4.1482,-2.3932) (4.1948,-2.5011) (4.2414,-2.6302) (4.2880,-2.7042) (4.3346,-3.0236) (4.3813,-3.3650) (4.4279,-3.2394) (4.4745,-2.9604) (4.5211,-2.1640) (4.5677,-2.5842) (4.6143,-2.2785) (4.6609,-2.3875) (4.7076,-2.4321) (4.7542,-4.2840) (4.8008,-2.6679) (4.8474,-3.3151) (4.8940,-3.2082) (4.9406,-4.8854) (4.9872,-3.0506) (5.0338,-2.8254) (5.0805,-2.2308) (5.1271,-2.5254) (5.1737,-2.3638) (5.2203,-2.3515) (5.2669,-3.2151) (5.3135,-2.6713) (5.3601,-2.8896) (5.4068,-3.2025) (5.4534,-3.7012) (5.5000,-3.0671) (5.5466,-2.2146) (5.5932,-2.6343) (5.6398,-2.3387) (5.6864,-2.4242) (5.7331,-2.5101) (5.7797,-2.7676) (5.8263,-2.2525) (5.8729,-3.0372) (5.9195,-2.1789) (5.9661,-3.2747) (6.0127,-2.3336) (6.0594,-2.7075) (6.1060,-2.6301) (6.1526,-2.4125) (6.1992,-2.3925) (6.2458,-2.6160) (6.2924,-2.2646) (6.3390,-3.0027) (6.3856,-2.6673) (6.4323,-3.3823) (6.4789,-3.3154) (6.5255,-2.4040) (6.5721,-2.6005) (6.6187,-2.4304) (6.6653,-2.4117)};
  \draw[black!45, dashed] (axis cs:0.9,-3) -- (axis cs:6.8,-3);
  \node[font=\tiny, text=black!55, anchor=south east] at (axis cs:6.7,-2.97) {``acceptable'' $=10^{-3}$};
\end{axis}
\end{tikzpicture}
\end{column}
\begin{column}{0.45\textwidth}
\small
The 45-degree test checked \emph{one point}. The residual checks \textbf{every} point, and it is less flattering:
\medskip
\begin{center}\footnotesize
\begin{tabular}{@{}lr@{}}
\toprule
mean $\log_{10}\mathcal{E}$ & $-2.48$ \\
worst $\log_{10}\mathcal{E}$ & $-1.81$ \\
\bottomrule
\end{tabular}
\end{center}
\medskip
A worst case of $10^{-1.8}$ is a \textbf{1.6\% consumption mistake} --- a decade worse than 3.A1's ``acceptable'' line.
\medskip
\begin{cautionbox}\footnotesize
And that is with $N=500$ and a converged value function. The staircase does not care how long you iterate.
\end{cautionbox}
\end{column}
\end{columns}
\end{frame}

%=============================================================================
\begin{frame}[fragile]{Version 1: Add a Shock, Change Three Lines}
\begin{columns}[T]
\begin{column}{0.54\textwidth}
\begin{lstlisting}
z  = np.array([0.95, 1.05])
Pi = np.array([[0.9, 0.1],
               [0.1, 0.9]])

# (1) cash on hand gains an axis
y = z[None,:]*(A*kg**alpha)[:,None] \
    + (1-delta)*kg[:,None]          # (N,2)
C = y[:,:,None] - kg[None,None,:]   # (N,2,N)

# (2) the expectation is a matrix product
EV = V @ Pi.T                       # (N,2)

# (3) one more axis to maximise over
M  = U + beta*EV.T[None,:,:]
Vn, pol = M.max(axis=2), M.argmax(axis=2)
\end{lstlisting}
\end{column}
\begin{column}{0.43\textwidth}
\small
\begin{center}\footnotesize
\begin{tabular}{@{}lr@{}}
\toprule
sweeps & 228 \\
wall time & 0.13\,s \\
return array & 4.0 MB \\
$g(k^{*})$, low $z$ & 3.7644 \\
$g(k^{*})$, high $z$ & 3.8793 \\
\bottomrule
\end{tabular}
\end{center}
\medskip
\begin{examplebox}\footnotesize
The two policies \textbf{bracket} the deterministic $k^{*}=3.8219$: a good shock saves more, a bad one saves less. A sign check you can do in your head, and it catches a transposed $\Pi$ immediately.
\end{examplebox}
\end{column}
\end{columns}
\vspace{1mm}
\src{The iteration count is unchanged --- it belongs to $\beta$, not to the state space (3.A2).}
\end{frame}

%=============================================================================
\begin{frame}{What We Deliberately Did Not Do}
\renewcommand{\arraystretch}{1.18}
\begin{center}\footnotesize
\begin{tabularx}{\linewidth}{@{}>{\raggedright\arraybackslash}p{3.1cm} >{\raggedright\arraybackslash}X c@{}}
\toprule
\textbf{In this code} & \textbf{What it costs, measured here} & \\
\midrule
$k'$ restricted to the grid & the crossing can only land within $h$ of $k^{*}$; worst residual $10^{-1.8}$ & S4 \\
No interpolation of $v$ & no way to evaluate a candidate between grid points & S4 \\
\texttt{argmax} over $N$ candidates & $57$ million comparisons for one deterministic solve & S4 \\
Every sweep a full sweep & 228 of them, when the policy settles far earlier & S4 \\
$\Pi$ written by hand & no connection to any estimated AR(1) & S4 \\
One continuous state & $(N,2,N)$ is 4 MB; a second state makes it 2 GB & S5 \\
\bottomrule
\end{tabularx}
\end{center}
\vspace{1mm}
\begin{takeawaybox}\small
Not one of these is a bug. They are the \textbf{cheapest possible choice} at each of the six decisions of 3.A1 --- and this deck is the measurement that says what each cheap choice cost. Next week we buy them back one at a time and re-run exactly this example.
\end{takeawaybox}
\end{frame}

%=============================================================================
\begin{frame}{Takeaways}
\begin{takeawaybox}
\begin{itemize}\small
  \item The whole solver is \textbf{five steps and about twenty-five lines}: parameters and grid, return matrix, Bellman step, loop, read the policy. Nothing in it is a numerical method.
  \item \textbf{Compute the benchmark before the solver.} Here it is one number --- $k^{*}=3.8219$ --- and one number is enough to catch most mistakes.
  \item The policy crossed the 45-degree line at \textbf{3.8161}, the nearest grid point to $k^{*}$. Right to within the grid, and no better: exactly what a correct implementation should give.
  \item One point passing is not the same as every point passing. The Euler residual over the whole grid is $10^{-1.8}$ at worst --- \textbf{a 1.6\% consumption mistake}, with a fully converged value function.
  \item Adding a shock changed \textbf{three lines} and the iteration count not at all. State space is memory; $\beta$ is time.
\end{itemize}
\end{takeawaybox}
\end{frame}

%=============================================================================
\begin{frame}{Bridge: Now the Language It Is Written In}
\begin{columns}[T]
\begin{column}{0.52\textwidth}
\small
Everything in this deck was NumPy: an array, a broadcast, a reduction along an axis, and one piece of fancy indexing. Five constructs, and they carried the entire algorithm.
\medskip

That is not a coincidence, and it is not something you can pick up by reading. Block B is those constructs from the ground up --- plus the environment they run in, and the one library that will hand you derivatives when the next method needs them.
\end{column}
\begin{column}{0.45\textwidth}
\begin{conceptbox}[title=\textbf{After the break (Block B)}]\small
\textbf{3.B1} Your toolkit: Python, environments, version control.\\[1mm]
\textbf{3.B2} From loops to arrays --- this deck's \texttt{U} and \texttt{M}, explained properly.\\[1mm]
\textbf{3.B3} PyTorch, and coding with AI.\\[1mm]
\textbf{3.B4} The data workflow.
\end{conceptbox}
\end{column}
\end{columns}
\end{frame}

\end{document}
